Title: More Structural Shocks than Observables in DSGE Estimation
Author: H. Hollander (notes compiled with AI assistance)
Purpose: Short technical note for reference and pedagogy (DSGE training)

1. State-space representation (why “counting shocks vs observables” is misleading)
A linear(ized) DSGE solution can be written in state-space form:

  Transition (states):
    x_t = F(θ) x_{t-1} + G(θ) ε_t,      ε_t ~ N(0, Q)

  Measurement (observables):
    y_t = H(θ) x_t + u_t,              u_t ~ N(0, R)

where x_t is a vector of latent states (not directly observed), y_t are observed data, and ε_t are structural shocks.
Shocks affect observables indirectly through the latent state vector and the model’s propagation mechanism.

Key references:
- An & Schorfheide (2007, JEDC)
- Herbst & Schorfheide (2015, book)
- Durbin & Koopman (2012, book)
- Smets & Wouters (2007, AER)

2. The Kalman filter and the likelihood (the key mathematical point)
The likelihood of the observables integrates out (marginalizes over) latent states:

  p(y_{1:T} | θ) = ∫ p(y_{1:T}, x_{1:T} | θ) dx_{1:T}
                 = ∫ [ Π_{t=1}^T p(y_t | x_t, θ) ] p(x_{1:T} | θ) dx_{1:T}

Under Gaussian assumptions, the Kalman filter computes this likelihood recursively via prediction errors:

  log L(θ) = Σ_{t=1}^T [ -½ log |S_t| - ½ ν_t' S_t^{-1} ν_t ] + constant

where:
- ν_t = y_t − E[y_t | y_{1:t−1}] is the one-step-ahead forecast error
- S_t = Var(ν_t | y_{1:t−1}) is the forecast error variance

Crucially, this does not require a one-to-one mapping between shocks and observables. Identification comes from the full
dynamic structure (propagation through F), cross-equation restrictions, and shock-specific loadings through G and H.

Key references:
- Durbin & Koopman (2012)
- Hamilton (1994, Ch. 13)
- An & Schorfheide (2007)

Useful explainers (accessible):
- Wikipedia (Kalman filter): https://en.wikipedia.org/wiki/Kalman_filter
- Durbin & Koopman (2012) is the standard state-space reference (book).
- An & Schorfheide (2007) is a widely-cited DSGE state-space primer in JEDC.

3. Why medium-scale NK models can sustain many shocks
In medium-scale NK DSGE models (e.g., Smets & Wouters 2007), shocks enter different structural blocks (Euler equations,
Phillips curves, capital accumulation, policy rules). Even if multiple shocks affect the same observable, they typically do so
with distinct persistence and propagation, producing distinguishable impulse responses (IRFs) and cross-covariances.

Reference:
- Smets & Wouters (2007, AER)

4. When “too many shocks” becomes a problem
Difficulties arise when:
(a) Shocks are observationally equivalent (similar equation entry + similar persistence + similar propagation), leading to weak identification.
(b) Shocks absorb structural misspecification, potentially distorting structural parameters and implied mechanisms.
(c) Many highly persistent shocks become collinear, increasing prior sensitivity and “variance swapping” across shocks.

References:
- Den Haan & Drechsel (2021, JME) on misspecification of structural disturbances and how flexible disturbances reveal it.
- Cúrdia & Reis (2012, JME) on correlated disturbances and implications for inference.

5. Identification criterion (what matters, formally)
The relevant condition is NOT:
  #shocks ≤ #observables

Instead, identification depends on rank / local invertibility conditions: whether the mapping from structural parameters
(and shock process parameters) into the likelihood-implied moments (or the state-space likelihood itself) is locally one-to-one.
Practically: distinct dynamic responses and cross-equation restrictions must pin down parameters and shock processes.

Reference:
- Iskrev (2010, JME) on local identification in DSGE models

6. Post-2008 financial-friction models (why shock counts rose, and what to watch)
Post-2008 DSGE models typically add financial wedges/shocks: risk premia, credit spreads, intermediation constraints,
capital quality/valuation, liquidity, and occasionally housing or collateral constraints. These shocks are often needed to match
spread dynamics, investment/credit comovement, and policy transmission. But added flexibility increases the risk that shocks
mimic each other or mask misspecification.

Good practice in these environments:
- Report variance decompositions and IRFs; check that key shocks imply economically interpretable dynamics.
- Use identification diagnostics (e.g., Iskrev-style / Dynare Identification Toolbox).
- Consider tests/robustness for shock specification and correlation (cf. Den Haan & Drechsel 2021; Cúrdia & Reis 2012).

Representative references:
- Bernanke, Gertler & Gilchrist (1999, QJE)
- Iacoviello (2005, AER)
- Gerali et al. (2010, JMCB)
- Gertler & Karadi (2011, JME)
- Christiano, Motto & Rostagno (2014, AER)

Summary
More shocks than observables is not inherently problematic in DSGE estimation. The Kalman filter computes the likelihood by
integrating over latent states, and identification comes from dynamic structure and cross-equation restrictions. Problems arise
when shocks are observationally equivalent or absorb misspecification—risks that are especially relevant in post-2008 macro-
financial models.